\section*{Bab \ref{ch:b1:radiation-heat-transport} --- Radiasi dan transpor panas}
% values: calc/ch27.py

\begin{solution}{exo:b1:radiation-heat-transport:1}
(a)~$S_0/\pi = 1361/\pi = \qty{433.2}{W.m^{-2}}$. (b)~Di khatulistiwa pada ekuinoks
$\bar Q = (S_0/\pi)(a/r)^2$, sehingga $(a/r)^2 = 436.7/433.2 = 1.0080$ dan $r/a =
1.0080^{-1/2} = 0.9960$: Bumi \qty{0.4}{\percent} lebih dekat ke Matahari daripada
jarak rata-ratanya. (c)~$\bar Q = 433.2 \times 0.9683 \times \cos\ang{23.44} =
\qty{384.9}{W.m^{-2}}$.
\end{solution}

\begin{solution}{exo:b1:radiation-heat-transport:2}
(a)~$415.6/172.4 = 2.41$. (b)~Zona di antara sejajar $\pm\varphi$ memiliki luas
$2\pi a_\oplus^2 \times 2\sin\varphi$, sebagian $\sin\varphi$ dari bola:
$\sin\ang{37.0} = 0.602$, \qty{60}{\percent}. (c)~Pada sumbu \omterm{def:b1:spherical-astronomy:geographic}{lintang}, pita dekat kutub
tampak selebar pita dekat khatulistiwa, tetapi meliputi luas yang jauh lebih kecil: zona
di sebelah kutub dari \ang{60} hanya $1 - \sin\ang{60} =
\qty{13.4}{\percent}$ dari permukaan. Terhadap $\sin\varphi$, tropik mengisi sebagian besar
sumbu.
\end{solution}

\begin{solution}{exo:b1:radiation-heat-transport:3}
(a)~$T_e = (\text{OLR}/\sigma)^{1/4}$: \qty{256.4}{K}, \qty{241.3}{K} dan
\qty{254.5}{K}. (b)~Sinar matahari terserap: $287.7/75.8 = 3.80$; OLR: $244.9/192.3 =
1.27$. (c)~\omterm{def:b1:radiation-heat-transport:net-radiation}{Radiasi bersih} mengikuti sinar matahari terserap: positif dan besar di tropik,
sangat negatif di \omterm{def:b1:spherical-astronomy:geographic}{lintang} tinggi. OLR, yang bervariasi sedikit, tidak mengimbangi
turunnya sinar matahari terserap.
\end{solution}

\begin{solution}{exo:b1:radiation-heat-transport:4}
(a)~$2.82 \times 365.25 = \qty{1030}{mm}$, \qty{1.03}{m} per tahun. (b)~$1.030 \times
78.1 = \qty{80.4}{W.m^{-2}}$.
(c)~$LE$ reanalisis \qty{2}{W.m^{-2}} lebih besar. Dalam banyak tahun atmosfer tidak dapat
terus mengumpulkan uap air: semua air yang menguap turun lagi, sehingga $LE$ dan $LP$
harus sama dalam rerata global. Selisihnya adalah ketidakseimbangan model, seperti
ketidakseimbangan radiatifnya.
\end{solution}

\begin{solution}{exo:b1:radiation-heat-transport:5}
(a)~$\bar Q = S_0(a/r)^2\sin\delta = 1361 \times 0.9683 \times 0.3978 =
\qty{524.2}{W.m^{-2}}$. (b)~$1361 \times 1.0334 \times 0.3978 =
\qty{559.5}{W.m^{-2}}$; rasionya $1.0334/0.9683 = 1.067$. (c)~Kutub memiliki
sinar matahari 24 jam dari 24 pada ketinggian tetap, sehingga $\sin h_\odot$ rerata atas
hari adalah $\sin\delta = 0.398$. Di khatulistiwa Matahari tinggi pada tengah hari tetapi
turun separuh waktu; rerata $\sin h_\odot$ atas 24 jam adalah $\cos\delta/\pi =
0.292$. Rasionya adalah $\pi\tan\delta = 1.362$
(\cref{prop:b1:radiation-heat-transport:polar-summer}).
\end{solution}

\begin{solution}{exo:b1:radiation-heat-transport:6}
(a)~$\cos h_0 = -\tan\ang{40}\tan\ang{23.44} = -0.3638$, $h_0 = \ang{111.33} =
\qty{1.943}{rad}$; hari berlangsung $24 \times 111.33/180 = \qty{14.84}{h}$; $\bar Q =
(1361/\pi) \times 0.9683 \times (1.943 \times 0.6428 \times 0.3978 + 0.7660 \times
0.9175 \times 0.9315) = \qty{483.0}{W.m^{-2}}$. (b)~$\cos h_0 = +0.3638$, $h_0 =
\ang{68.67}$, \qty{9.16}{h}; $\bar Q = \qty{155.9}{W.m^{-2}}$. (c)~Rasionya adalah
$483.0/155.9 = 3.10$. Jarak berkontribusi $0.9683/1.0334 = 0.937$: ia bekerja melawan
musim panas belahan utara. Geometri, ketinggian Matahari dan panjang hari,
berkontribusi $3.10/0.937 = 3.31$, seluruh kontras musiman dan lebih.
\end{solution}

\begin{solution}{exo:b1:radiation-heat-transport:7}
(a)~$B = H/LE$: \num{0.093} di khatulistiwa, \num{0.374} pada \ang{25} LU,
\num{0.338} pada \ang{45} LU, \num{-0.152} pada \ang{70} LU. (b)~\ang{25} LU: rerata
zonal mencakup gurun panas yang besar, daratan kering dengan sedikit air untuk diuapkan, tempat
tanah yang panas memanaskan udara. (c)~$H < 0$ berarti panas mengalir dari udara ke
permukaan: dalam rerata tahunan salju dan es laut lebih dingin daripada udara di atasnya,
yang dipanaskan oleh udara yang dibawa dari \omterm{def:b1:spherical-astronomy:geographic}{lintang} yang lebih rendah.
\end{solution}

\begin{solution}{exo:b1:radiation-heat-transport:8}
(a)~Dinilai: $\alpha = 100/340 = 0.294$, $R = 340 - 100 - 239 =
\qty{+1}{W.m^{-2}}$ (\num{0.7} dibulatkan). Reanalisis: $\alpha = 117.0/341.9 =
0.342$, $R = 341.9 - 117.0 - 237.7 = \qty{-12.8}{W.m^{-2}}$ dengan suku-suku yang
dibulatkan ini, \qty{-12.9}{W.m^{-2}} dengan suku-suku yang tidak dibulatkan dari
\cref{ex:b1:radiation-heat-transport:global}. (b)~Reanalisis menyerap
$224.8$ dibanding $240$, \qty{15.2}{W.m^{-2}} lebih sedikit, karena ia memantulkan
\qty{17}{W.m^{-2}} lebih banyak; OLR-nya hanya \qty{1.3}{W.m^{-2}} lebih rendah. Pantulan
berlebih menjelaskan ketidakseimbangan itu. (c)~Mengintegrasikan $R$ dari Kutub Selatan ke
Kutub Utara memberi integral global, $-12.9 \times 4\pi a_\oplus^2 = -12.9 \times
5.10\times10^{14} = \qty{-6.6}{PW}$. Transpor melintasi Kutub Utara tidak mungkin: tidak ada
apa pun di luarnya, sejajar itu adalah sebuah titik. Proposisi hanya berlaku bila
$\int R\,\mathrm dA = 0$, yang diwajibkan oleh keadaan tunak.
\end{solution}

\begin{solution}{exo:b1:radiation-heat-transport:9}
(a)~Luasnya $2\pi a_\oplus^2 \times 2\sin\ang{30} = 2\pi a_\oplus^2 =
\qty{2.55e14}{m^2}$, setengah Bumi; surplusnya $33 \times 2.55\times10^{14} =
\qty{8.4}{PW}$. (b)~\qty{4.2}{PW} melintasi setiap sejajar. (c)~Keduanya berjumlah
\qty{8.5}{PW}, surplusnya, tetapi tidak sama: belahan selatan mengekspor
\qty{0.5}{PW} lebih banyak daripada belahan utara. Surplusnya tidak dibagi di
khatulistiwa: dalam reanalisis ini \qty{0.36}{PW} dari surplus tropik utara melintasi
khatulistiwa ke selatan.
\end{solution}

\begin{solution}{exo:b1:radiation-heat-transport:10}
(a)~Tulis rerata harian sebagai $\bar Q = (a/r)^2 F(\varphi, \delta)$, dengan $F$ kurung
\cref{thm:b1:radiation-heat-transport:daily} dikali $S_0/\pi$. Selama $\mathrm dt$ energi
yang diterima per meter persegi adalah $\bar Q\,\mathrm dt$, dan
$\mathrm dt = r^2\,\mathrm d\lambda/h$, sehingga $\bar Q\,\mathrm dt = (a^2/h)
F(\varphi, \delta(\lambda))\,\mathrm d\lambda$: faktor $r^2$ saling meniadakan.
(b)~Selama setahun energinya adalah $(a^2/h)\int_0^{2\pi} F(\varphi,
\delta(\lambda))\,\mathrm d\lambda$. Rumusnya tidak berubah bila $\varphi$ dan
$\delta$ berubah tanda bersama-sama, $F(-\varphi, \delta) = F(\varphi, -\delta)$, dan
$\delta(\lambda + \pi) = -\delta(\lambda)$, karena $\sin\delta = \sin\varepsilon
\sin\lambda$. Menggeser $\lambda$ sebesar $\pi$ dalam integral menunjukkan bahwa $-\varphi$
menerima energi yang sama dengan $\varphi$. (c)~Musim semi dan panas belahan utara
berlangsung \num{7.6} hari lebih lama, karena Bumi dekat aphelium dan bergerak lambat saat
itu; musim panas di selatan lebih terang tetapi lebih pendek. Dengan (a), kedua efek
itu tepat saling meniadakan.
\end{solution}

\begin{solution}{exo:b1:radiation-heat-transport:11}
(a)~$\langle (a/r)^2\rangle = \frac1P\int_0^P \frac{a^2}{r^2}\,\mathrm dt =
\frac1P\int_0^{2\pi}\frac{a^2}{r^2}\frac{r^2}{h}\,\mathrm d\lambda = \frac{2\pi
a^2}{hP} = \frac{2\pi a^2}{2\pi a b} = \frac ab$, dan $b = a\sqrt{1-e^2}$.
(b)~Rerata sesaat dari global $\frac14 S_0(a/r)^2$ menjadi $\frac14 S_0
\langle(a/r)^2\rangle = S_0/(4\sqrt{1-e^2})$. (c)~$340.25/\sqrt{1 - 0.0167^2} =
\qty{340.30}{W.m^{-2}}$, \qty{0.05}{W.m^{-2}} atau \qty{0.014}{\percent} di atas
$S_0/4$.
\end{solution}

\begin{solution}{exo:b1:radiation-heat-transport:12}
(a)~Atmosfer $0.78 \times 4.10 = \qty{3.20}{PW}$, samudra $0.22 \times 4.10 =
\qty{0.90}{PW}$. (b)~Atmosfer saja ditemukan membawa lebih banyak, \qty{5.0}{PW},
daripada atmosfer dan samudra reanalisis bersama, \qty{4.02}{PW}: transpor total
reanalisis kurang setidaknya seperlima ($1 - 4.02/5.0 = 0.20$).
\omterm{def:b1:radiation-heat-transport:net-radiation}{Radiasi bersihnya} bervariasi terlalu sedikit antara tropik dan kutub. (c)~Mengurangkan
konstanta dari $R$ membuat integral globalnya nol tetapi tidak mengubah variasinya terhadap
\omterm{def:b1:spherical-astronomy:geographic}{lintang}; galatnya terletak pada variasi itu, pada albedo awan dan permukaan model di setiap
\omterm{def:b1:spherical-astronomy:geographic}{lintang}, bukan hanya pada rerata global.
\end{solution}

\begin{solution}{pb:b1:radiation-heat-transport:1}
\begin{enumerate}
  \item Juni: $\cos h_0 = -0.3638$, $h_0 = \ang{111.33}$, $\bar Q =
        \qty{483.0}{W.m^{-2}}$. Desember: $h_0 = \ang{68.67}$, $\bar Q =
        \qty{155.9}{W.m^{-2}}$. Maret: $h_0 = \ang{90}$, $\bar Q = (1361/\pi)
        \times 1.0081 \times \cos\ang{40} = \qty{334.6}{W.m^{-2}}$.
  \item Juni: karena $\tan\ang{80}\tan\ang{23.44} = 2.46 > 1$, siang kutub, $h_0 = \pi$:
        $\bar Q = 1361 \times 0.9683 \times \sin\ang{80}\sin\ang{23.44} =
        \qty{516.3}{W.m^{-2}}$. Desember: malam kutub, $\bar Q = 0$. Maret:
        $(1361/\pi) \times 1.0081 \times \cos\ang{80} = \qty{75.8}{W.m^{-2}}$.
  \item Pada \ang{80} LU Matahari tidak terbenam ($h_0 = \ang{180}$, 24 jam) sedangkan pada
        \ang{40} LU ia bersinar selama \qty{14.8}{h}. Ketinggiannya lebih rendah, antara
        \ang{13.4} dan \ang{33.4}, tetapi suku $h_0\sin\varphi\sin\delta$, yang
        bertambah dengan panjang hari dan \omterm{def:b1:spherical-astronomy:geographic}{lintang}, menang.
  \item Nilai pita melebihi nilai pusat sebesar \qty{0.5}{\percent},
        \qty{0.8}{\percent}, \qty{1.7}{\percent}, \qty{1.7}{\percent} dan
        \qty{1.3}{\percent}. Pertama, reanalisis memakai iradiasi matahari yang sedikit lebih besar
        (Pertanyaan~5), \qty{0.5}{\percent} lebih banyak. Kedua, rerata pita bukan nilai di pusat:
        setengah pita yang menuju khatulistiwa memiliki luas dan sinar matahari lebih banyak daripada
        setengah yang menuju kutub.
  \item Dengan \cref{prop:b1:radiation-heat-transport:global}, $S_0 = 4\sqrt{1-e^2}
        \times 341.9 = \qty{1367}{W.m^{-2}}$, \qty{0.47}{\percent} lebih banyak daripada \qty{1361}{W.m^{-2}}.
  \item $A = 2\pi a_\oplus^2(1 - \sin\ang{40}) = 2\pi \times (6.371\times10^6)^2
        \times 0.3572 = \qty{9.110e13}{m^2}$, sebagian $(1 - \sin\ang{40})/2 =
        0.179$ dari Bumi.
  \item $(1 - \alpha)\bar Q$: \num{198.0}, \num{159.3}, \num{124.7}, \num{92.5},
        \num{74.7} \unit{W.m^{-2}}.
  \item $R$: \num{-28.1}, \num{-51.0}, \num{-76.6}, \num{-104.9}, \num{-116.9}
        \unit{W.m^{-2}}. Kelima pita berada dalam defisit; pita kutub
        \ang{80}--\ang{90} memiliki yang terbesar.
  \item $T_e = (\text{OLR}/\sigma)^{1/4}$: \num{251.3}, \num{246.8}, \num{244.1},
        \num{242.9}, \num{241.1} \unit{K}.
  \item Terserap: $198.0/74.7 = 2.65$; OLR: $226.1/191.6 = 1.18$.
  \item Matahari lebih rendah, dan cahaya yang datang dengan sudut miring lebih banyak
        dipantulkan oleh awan dan air; salju dan es laut, permukaan cerah, menutupi lebih banyak
        tanah dan laut ke arah kutub.
  \item OLR sebanding dengan $T_e^4$, dan $T_e$ bervariasi hanya \qty{10}{K}, \qty{4}{\percent},
        di seluruh tudung: \omterm{def:b1:spherical-astronomy:geographic}{lintang} tinggi dijaga jauh lebih hangat daripada yang diizinkan
        sinar mataharinya saja, oleh panas yang dibawa dari selatan.
  \item Dalam \unit{m^2}: \num{3.143e13}, \num{2.550e13}, \num{1.879e13},
        \num{1.151e13}, \num{0.387e13}; jumlahnya adalah \qty{9.110e13}{m^2}.
  \item $R_iA_i$: \num{-0.883}, \num{-1.300}, \num{-1.439}, \num{-1.207},
        \num{-0.453} \unit{PW}; tudung: \qty{-5.28}{PW}.
  \item $\alpha_\text{tudung} = \sum\alpha_i\bar Q_iA_i/\sum\bar Q_iA_i = 0.400$,
        dibandingkan dengan \num{0.342} untuk globe.
  \item $-12.90 \times 5.10\times10^{14} = \qty{-6.58}{PW}$: Bumi akan kehilangan panas
        dengan laju itu dan mendingin. Ketidakseimbangan yang dinilai adalah
        $+0.7 \times 5.10\times10^{14} = \qty{+0.36}{PW}$, delapan belas kali lebih kecil dan
        tandanya berlawanan: ketidakseimbangan reanalisis adalah galat model.
  \item $R_i + 12.90$ terkoreksi: \num{-15.2}, \num{-38.1}, \num{-63.7},
        \num{-92.0}, \num{-104.0} \unit{W.m^{-2}}; daya \num{-0.478},
        \num{-0.972}, \num{-1.197}, \num{-1.059}, \num{-0.403} \unit{PW}; tudung:
        \qty{-4.108}{PW}, defisit rerata \qty{45.1}{W.m^{-2}}.
  \item Energi tudung tidak berubah dari tahun ke tahun, sehingga apa yang hilang di
        puncak atmosfer harus masuk ke dalamnya, dan satu-satunya jalan masuk adalah melintasi
        sejajar ke-40: $T(\ang{40}) = -2\pi a_\oplus^2\int_{\ang{40}}^{\ang{90}}
        R\cos\varphi\,\mathrm d\varphi$.
  \item $T(\ang{40}\,\text{LU}) = +\qty{4.11}{PW}$, ke utara.
  \item $LP = 2.4655\times10^6 \times 1.95/86400 = \qty{55.6}{W.m^{-2}}$, lebih banyak
        daripada $LE = \qty{42.4}{W.m^{-2}}$: lebih banyak air jatuh di tudung daripada yang
        menguap darinya; selisihnya tiba sebagai uap.
  \item $(LE - LP) - 2.07 = 42.4 - 55.6 - 2.07 = \qty{-15.3}{W.m^{-2}}$; dikali
        luas tudung: \qty{-1.40}{PW}. Atmosfer membawa \qty{1.40}{PW}
        panas laten melintasi \ang{40} LU.
  \item $1.40/4.11 = 0.34$: sepertiga.
  \item Samudra $0.22 \times 4.11 = \qty{0.90}{PW}$; atmosfer $0.78 \times 4.11 =
        \qty{3.20}{PW}$, dari jumlah itu \qty{1.40}{PW} laten dan \qty{1.80}{PW}
        dibawa sebagai panas sensibel dan energi potensial udara.
  \item Batas bawah. Total reanalisis pada \ang{43} LU, \qty{4.02}{PW}, lebih kecil
        daripada atmosfer saja dalam studi, \qty{5.0}{PW}; total yang sebenarnya
        lebih besar.
  \item Sekitar \qty{4.11}{PW}, \qty{4.11e15}{W}, menurut radiasi reanalisis, sepertiganya
        sebagai panas laten; nilai yang sebenarnya masih lebih besar.
\end{enumerate}
\end{solution}
